\documentclass[10pt]{beamer}

\geometry{paperwidth=10in,paperheight=5.625in,margin=0in}
\usepackage{beamerthemeSmartWaterTU}
\usepackage{amsmath}
\usepackage{booktabs}
\usepackage{array}

\title{Can a GAN Detect and Localize Leaks from Pressure Images?}
\author{Cemal Hekim, Jan Lcukey, and Moyang Chen}
\institute{Rajabi, Komeilian, Wan, and Farmani - Water Research 238 (2023) 120012}

\newcommand{\PaperAsset}[2][]{\includegraphics[#1]{paper-assets/#2}}
\newcommand{\PaperSource}[1]{%
  \vspace{0.035in}{\fontsize{8}{9}\selectfont\color{black!65}\raggedleft
  Source: Rajabi et al. (2023), Water Research 238, 120012, #1.\par}}
\newcommand{\LitSource}[1]{%
  \vspace{0.025in}{\fontsize{8}{9}\selectfont\color{black!65}\raggedleft #1\par}}

\newcommand{\JCBullets}[1]{%
  \begin{itemize}
    \setlength{\itemsep}{0.08in}
    \setlength{\parskip}{0pt}
    #1
  \end{itemize}
}

\newcommand{\DenseBullets}[1]{%
  \begin{itemize}
    \setlength{\itemsep}{0.025in}
    \setlength{\parskip}{0pt}
    #1
  \end{itemize}
}

\newcommand{\CompactBullets}[1]{%
  \begin{itemize}
    \setlength{\itemsep}{0.012in}
    \setlength{\parskip}{0pt}
    #1
  \end{itemize}
}

\begin{document}

% MAIN TALK: 11:30 target, leaving ~30 s safety margin.
% Presenter balance:
% Cemal: slides 1--4 (3:50)
% Jan: slides 5--8 (3:50)
% Moyang: slides 9--13 (3:50)

\SWNTitleSlide
  {Can a GAN Detect and Localize Leaks from Pressure Images?}
  {Cemal Hekim, Jan Lcukey, and Moyang Chen}
  {Rajabi et al., Water Research 238 (2023), 120012}
  { -- critical paper presentation}

% 1:00 -- Cemal
\begin{SWNFrame}{RESEARCH CHALLENGES}
\begin{columns}[T,onlytextwidth]
\begin{column}{0.51\textwidth}
\textbf{Challenge 1: Can weak leak signals be separated from noise?}
\JCBullets{
  \item Pressure sensors are affordable and real-time, but distant or small leaks may produce weak signals.
  \item Demand uncertainty, sensor error, and operational changes can resemble a leak.
}
\end{column}
\begin{column}{0.45\textwidth}
\textbf{Challenge 2: Can detection and localization work without labeled leaks?}
\JCBullets{
  \item Labeled field leaks are scarce; exhaustive leak simulation is expensive.
  \item Existing data-driven methods struggle with scarce labels and spatial localization.
  \item Model-based methods require calibration and repeated hydraulic simulation.
}
\end{column}
\end{columns}
\vspace{0.04in}
\textbf{Research need:} a fast, uncertainty-aware, spatial method that learns without labeled leaks and narrows the field inspection area.
\LitSource{Context synthesized from Rajabi et al. (2023), pp. 1--3; Wan et al. (2022).}
\end{SWNFrame}

% 1:15 -- Cemal
\begin{SWNFrame}{PAPER GOAL AND RESEARCH QUESTION}
\begin{columns}[T,onlytextwidth]
\begin{column}{0.57\textwidth}
\textbf{Research question derived from the paper's objective}
\vspace{0.06in}

\emph{Can a conditional deep convolutional GAN, trained only on leak-free
demand--pressure behavior, detect and spatially localize unseen leaks from
pressure observations in near real time?}

\vspace{0.12in}
\textbf{Success would require}
\DenseBullets{
  \item accurate detection despite demand and measurement uncertainty;
  \item useful spatial localization from pressure observations;
  \item low-cost inference suitable for real-time operation.
}
\end{column}
\begin{column}{0.39\textwidth}
\textbf{Goal set by the authors}
\JCBullets{
  \item Develop a robust methodology for accurate, real-time leak detection and localization in complex settings.
  \item Address uncertainty, scale, scarce labels, and complex features while using pressure data.
}
\vspace{0.08in}
\textbf{Not the goal:} ``use a GAN.'' The GAN is the authors' proposed means of achieving the goal.
\end{column}
\end{columns}
\PaperSource{pp. 2--3}
\end{SWNFrame}

% 1:15 -- Cemal
\begin{SWNFrame}{THE CORE IDEA IN ONE PIPELINE}
\begin{columns}[T,onlytextwidth]
\begin{column}{0.58\textwidth}
\PaperAsset[width=\linewidth]{fig1-flowchart.png}
\PaperSource{Fig. 1, p. 3}
\end{column}
\begin{column}{0.38\textwidth}
\textbf{Four steps}
\JCBullets{
  \item Simulate leak-free demand and pressure with EPANET.
  \item Convert point values into 256 x 256 spatial images using kriging.
  \item Learn the normal mapping: demand image $\rightarrow$ pressure image.
  \item Compare predicted and observed pressure images using SSIM.
}
\vspace{0.04in}
\textbf{Novel combination:} hydraulic simulation supplies normal data; the neural network becomes a fast surrogate; SSIM supplies the anomaly signal.
\end{column}
\end{columns}
\end{SWNFrame}

% 0:55 -- Jan
\begin{SWNFrame}{HOW THE METHOD MAKES A DECISION}
\begin{columns}[T,onlytextwidth]
\begin{column}{0.53\textwidth}
\PaperAsset[width=\linewidth]{fig2-architecture.png}
\PaperSource{Fig. 2, p. 5}
\end{column}
\begin{column}{0.43\textwidth}
\textbf{Prediction model}
\DenseBullets{
  \item Pix2pix-style conditional GAN.
  \item Modified U-Net generator with skip connections.
  \item PatchGAN discriminator; adversarial + L1 loss.
}
\textbf{Detection and localization}
\DenseBullets{
  \item Global $SSIM_o$: alarm when below an hourly mean $-3\sigma$ threshold for five consecutive steps.
  \item Local $SSIM_l$: minimum-similarity region estimates the leak location.
}
\textbf{Built-in tradeoff:} persistence suppresses false alarms but delays detection of small or short leaks.
\end{column}
\end{columns}
\end{SWNFrame}

% 1:00 -- Jan
\begin{SWNFrame}{EVALUATION DESIGN}
\begin{columns}[T,onlytextwidth]
\begin{column}{0.39\textwidth}
\PaperAsset[width=\linewidth]{fig3-map.png}
\PaperSource{Fig. 3, p. 6}
\end{column}
\begin{column}{0.57\textwidth}
\textbf{L-Town benchmark}
\DenseBullets{
  \item 782 junctions, 905 pipes, four pressure zones, and 33 sensors.
  \item One year of 5-minute data: 105,120 leak-free image pairs.
  \item Two image models: 780 observation points vs. 33 sensor points.
}
\textbf{Evaluation sets}
\DenseBullets{
  \item \textbf{Controlled:} 24 isolated, step-growing leaks at eight locations and three maximum leak areas.
  \item \textbf{More realistic:} six non-concurrent BattLeDIM leaks across one year.
}
\textbf{Metrics:} TPR, TNR, accuracy, detection time, and graph distance to the true leak.
\end{column}
\end{columns}
\end{SWNFrame}

% 1:00 -- Jan
\begin{SWNFrame}{RESULT 1: MAGNITUDE DRIVES DETECTION}
\begin{columns}[T,onlytextwidth]
\begin{column}{0.56\textwidth}
\centering
\PaperAsset[height=2.18in,keepaspectratio]{fig10-ssimo.png}
\vspace{0.025in}

\PaperAsset[height=0.62in,keepaspectratio]{fig11-local-ssim.png}
\PaperSource{Figs. 10--11, pp. 9--11}
\end{column}
\begin{column}{0.40\textwidth}
\textbf{At one example location}
\JCBullets{
  \item Large leak: detected almost immediately; localization converges to 4 m.
  \item Medium leak: detected after about 16 h and only after sufficient growth.
  \item Small leak: detected after about 21 h, near its maximum rate; localization about 96 m.
}
\textbf{Interpretation:} the method detects a sufficiently strong pressure-image disturbance, not the mere onset of leakage.
\end{column}
\end{columns}
\end{SWNFrame}

% 0:55 -- Jan
\begin{SWNFrame}{RESULT 2: AVERAGES HIDE FAILURES}
\begin{columns}[T,onlytextwidth]
\begin{column}{0.55\textwidth}
\PaperAsset[width=0.50\linewidth]{fig13-tpr-tnr.png}
\hfill
\PaperAsset[width=0.47\linewidth]{fig15-gdrl-dt.png}
\PaperSource{Figs. 13 and 15, pp. 10--13}
\end{column}
\begin{column}{0.41\textwidth}
\textbf{Reported evidence}
\DenseBullets{
  \item Average accuracy: 0.91 with 780 points; 0.90 with 33 sensors.
  \item TNR is generally high.
  \item TPR approaches 1 for the largest leaks, but approaches 0 for the smallest.
  \item Larger leaks are detected faster and localized more accurately.
}
\textbf{Our reading:} similar average accuracy does \emph{not} establish sensor-count invariance; the class balance and high TNR can mask missed small leaks.
\end{column}
\end{columns}
\end{SWNFrame}

% 0:50 -- Moyang
\begin{SWNFrame}{RESULT 3: YEAR-LONG STRESS TEST}
\begin{columns}[T,onlytextwidth]
\begin{column}{0.60\textwidth}
\PaperAsset[height=2.70in,keepaspectratio]{fig16-realtime.png}
\PaperSource{Fig. 16 and Table 2, pp. 12--13}
\end{column}
\begin{column}{0.36\textwidth}
\textbf{Six BattLeDIM leaks}
\DenseBullets{
  \item Five detected; the shortest abrupt event was missed.
  \item Detection generally occurred near peak leak rate.
  \item Whole-year accuracy: about 70\%.
  \item Localization error for detected leaks: 160--185 m.
}
\textbf{Meaning:} feasibility is shown, but early warning and field-ready localization are not. This is the strongest test because the leaks are smaller and time-varying.
\end{column}
\end{columns}
\end{SWNFrame}

% 0:55 -- Moyang
\begin{SWNFrame}{CRITICAL EVALUATION}
\begin{columns}[T,onlytextwidth]
\begin{column}{0.48\textwidth}
\textbf{Supported by the experiments}
\JCBullets{
  \item The pipeline can detect and approximately localize sufficiently large, isolated leaks.
  \item Inference can replace repeated online hydraulic simulations.
  \item Sparse 33-sensor images retained similar aggregate performance in the controlled scenarios.
}
\end{column}
\begin{column}{0.48\textwidth}
\textbf{Not yet established}
\JCBullets{
  \item That the GAN is better than a deterministic U-Net or residual baseline.
  \item Robustness to sensor faults, missing data, valve/pump changes, network drift, or simultaneous leaks.
  \item Generalization beyond one simulated network and its kriging/scaling choices.
}
\end{column}
\end{columns}
\vspace{0.04in}
\textbf{Validity judgment:} the conclusions are credible as a proof of concept, but the authors' robustness and practical-readiness claims are broader than the presented evidence.
\LitSource{Our evaluation, based on the experiments and assumptions reported by Rajabi et al. (2023), pp. 6--13.}
\end{SWNFrame}

% 0:45 -- Moyang
\begin{SWNFrame}{ASSUMPTIONS THAT MATTER IN PRACTICE}
\begin{columns}[T,onlytextwidth]
\begin{column}{0.49\textwidth}
\textbf{Model and data assumptions}
\DenseBullets{
  \item Known and sufficiently calibrated network.
  \item No initial defects; fixed topology and pipe characteristics.
  \item Normally distributed demand uncertainty and measurement error.
  \item Accurate, synchronized pressure data.
  \item No concurrent leaks in the year-long test.
}
\end{column}
\begin{column}{0.47\textwidth}
\textbf{Consequences if violated}
\DenseBullets{
  \item Normal operational change may trigger false alarms.
  \item Domain shift can corrupt the learned demand--pressure mapping.
  \item Kriging may create spatial patterns not supported by sparse measurements.
  \item A fixed $3\sigma$ threshold may not match utility-specific alarm costs.
}
\end{column}
\end{columns}
\vspace{0.05in}
\textbf{Key distinction:} the first column is reported by the authors; the second is our interpretation of deployment risk.
\PaperSource{limitations and future work, p. 13}
\end{SWNFrame}

% 0:50 -- Moyang
\begin{SWNFrame}{CONCLUSION, IMPLICATIONS, AND NEXT TESTS}
\begin{columns}[T,onlytextwidth]
\begin{column}{0.48\textwidth}
\textbf{Answer to the research question}
\JCBullets{
  \item \textbf{Yes, conditionally:} the method detects and localizes leaks once their pressure signature is strong enough.
  \item \textbf{No, not yet generally:} early small-leak detection and precise field localization remain weak.
}
\textbf{Practical implication}

Best viewed as a screening tool that prioritizes inspection areas, not as an autonomous leak pinpointing system.
\end{column}
\begin{column}{0.48\textwidth}
\textbf{Specific next experiments}
\DenseBullets{
  \item Compare CDCGAN against U-Net, persistence, and hydraulic-residual baselines.
  \item Test missing/faulty sensors, topology changes, drift, and simultaneous leaks.
  \item Validate on multiple networks and real utility measurements.
  \item Calibrate thresholds to detection delay and false-alarm cost.
  \item Report uncertainty for alarms and locations.
}
\end{column}
\end{columns}
\PaperSource{authors' conclusions plus our proposed evaluation priorities, p. 13}
\end{SWNFrame}

% 0:30 -- all presenters
\begin{SWNFrame}{TAKE-HOME VERDICT}
\begin{columns}[T,onlytextwidth]
\begin{column}{0.31\textwidth}
\textbf{Innovation}

Normal hydraulic behavior is learned as an image-to-image mapping; global and local SSIM turn prediction error into detection and localization.
\end{column}
\begin{column}{0.31\textwidth}
\textbf{Best evidence}

Large controlled leaks are detected quickly, while the year-long BattLeDIM test detects five of six events with about 70\% accuracy.
\end{column}
\begin{column}{0.31\textwidth}
\textbf{Our verdict}

Promising proof of concept, but comparative baselines and field robustness tests are needed before deployment claims are justified.
\end{column}
\end{columns}
\vspace{0.23in}
\centering
\textbf{Questions?}
\end{SWNFrame}

% BACKUP SLIDES -- open only during Q&A.
\appendix

\begin{SWNFrame}{BACKUP: WHY A GAN?}
\begin{columns}[T,onlytextwidth]
\begin{column}{0.48\textwidth}
\textbf{Potential GAN advantage}
\JCBullets{
  \item Adversarial training may preserve realistic spatial texture and sharper local structure.
  \item This could make SSIM residual maps more spatially informative.
}
\end{column}
\begin{column}{0.48\textwidth}
\textbf{What is missing}
\JCBullets{
  \item No deterministic U-Net ablation.
  \item No comparison of detection metrics, localization error, stability, or training cost.
  \item Therefore, the incremental value of the discriminator is unknown.
}
\end{column}
\end{columns}
\vspace{0.10in}
\textbf{Concise Q\&A answer:} the architecture is plausible, but the paper demonstrates the pipeline's performance, not the necessity of the GAN component.
\end{SWNFrame}

\begin{SWNFrame}{BACKUP: 33 VS. 780 POINTS}
\begin{columns}[T,onlytextwidth]
\begin{column}{0.52\textwidth}
\PaperAsset[width=\linewidth]{fig4-maps.png}
\PaperSource{Fig. 4, p. 7}
\end{column}
\begin{column}{0.44\textwidth}
\textbf{Reported result}
\DenseBullets{
  \item Average controlled-scenario accuracy differs by only 0.01.
  \item Pressure prediction MAPE is lower for the 33-point model.
}
\textbf{Why caution is needed}
\DenseBullets{
  \item Both are interpolated images, not independent spatial measurements.
  \item Separate scaling by pressure zone changes the learning problem.
  \item MAPE and aggregate accuracy do not test early small-leak sensitivity.
}
\textbf{Needed test:} repeat sensor subsets and placements; report uncertainty.
\end{column}
\end{columns}
\end{SWNFrame}

\begin{SWNFrame}{BACKUP: INTERPRETING 70\% ACCURACY}
\begin{columns}[T,onlytextwidth]
\begin{column}{0.49\textwidth}
\textbf{Why accuracy can mislead}
\DenseBullets{
  \item A one-year time series contains many more normal than leaking time steps.
  \item High specificity can dominate accuracy even when early leak sensitivity is poor.
  \item Five detected events do not describe alarm frequency or event-level precision.
}
\end{column}
\begin{column}{0.47\textwidth}
\textbf{Better reporting}
\DenseBullets{
  \item Event-level precision, recall, and false alarms per month.
  \item Detection delay from onset and from economically relevant leak rate.
  \item Precision--recall curve across thresholds.
  \item Localization radius at fixed confidence.
}
\end{column}
\end{columns}
\vspace{0.10in}
\textbf{Concise Q\&A answer:} 70\% supports partial feasibility, but it is insufficient by itself to establish operational usefulness.
\end{SWNFrame}

\begin{SWNFrame}{BACKUP: THRESHOLD AND FAILURE MODES}
\begin{columns}[T,onlytextwidth]
\begin{column}{0.48\textwidth}
\textbf{Decision rule}
\[
SSIM_o < \mu_{\mathrm{normal,hour}}-3\sigma_{\mathrm{normal,hour}}
\]
for five consecutive 5-minute steps.

\textbf{Benefit:} fewer transient false alarms.
\end{column}
\begin{column}{0.48\textwidth}
\textbf{Likely failure modes}
\DenseBullets{
  \item Short leak ends before five-step confirmation.
  \item Slow leak is detected only near peak magnitude.
  \item Non-stationary operations shift the normal distribution.
  \item Correlated or heavy-tailed errors violate the $3\sigma$ intuition.
}
\end{column}
\end{columns}
\vspace{0.10in}
\textbf{Improvement:} tune a cost-sensitive sequential detector on utility-specific false-alarm and delay requirements.
\end{SWNFrame}

\begin{SWNFrame}{BACKUP: CORE REFERENCES}
\fontsize{12}{15}\selectfont
\begin{thebibliography}{9}
\bibitem{rajabi}
Rajabi, M. M., Komeilian, P., Wan, X., \& Farmani, R. (2023).
Leak detection and localization in water distribution networks using conditional
deep convolutional generative adversarial networks.
\emph{Water Research, 238}, 120012.

\bibitem{wan}
Wan, X., Kuhanestani, P. K., Farmani, R., \& Keedwell, E. (2022).
Literature review of data analytics for leak detection in water distribution
networks: a focus on pressure and flow smart sensors.
\emph{Journal of Water Resources Planning and Management, 148}(10), 03122002.

\bibitem{isola}
Isola, P., Zhu, J.-Y., Zhou, T., \& Efros, A. A. (2017).
Image-to-image translation with conditional adversarial networks.
\emph{CVPR}, 1125--1134.

\bibitem{vrachimis}
Vrachimis, S. G. et al. (2022).
Battle of the leakage detection and isolation methods.
\emph{Journal of Water Resources Planning and Management, 148}(12), 04022068.
\end{thebibliography}
\end{SWNFrame}

\end{document}
